% chap02.tex 第二章 控制系统的数学模型（大气排版版）
\chapter{控制系统的数学模型}

% ===== 局部排版优化（仅影响本章，使版面更舒展）=====
\setlength{\parskip}{0.8ex plus 0.3ex minus 0.2ex}   % 段落间距增加
\setlist{
    leftmargin=2.5em,          % 列表左缩进增大
    itemsep=0.6em,             % 条目垂直间距
    parsep=0.3em,              % 条目内段落间距
    topsep=0.8em,              % 列表与上下文间距
    partopsep=0.3em
}
\renewcommand{\arraystretch}{1.5}      % 表格行高增加
\setlength{\tabcolsep}{10pt}           % 表格列间距增加
\addtolength{\abovedisplayskip}{2pt}   % 公式前间距微增
\addtolength{\belowdisplayskip}{2pt}   % 公式后间距微增
% ========================================

\bigskip  % 章节标题后额外留白

\section{本章考核要点}
\begin{enumerate}
    \item 依据物理定律列写系统微分方程的通用步骤；
    \item 机电典型元件微分方程、机电相似性（RC电路$\leftrightarrow$弹簧阻尼质量）；
    \item 微分方程$\leftrightarrow$传递函数$\leftrightarrow$频率特性三者变换关系；
    \item 方块图化简、信号流图+梅森增益公式求总传递函数；
    \item 小信号法非线性系统局部线性化；
    \item 脉冲传递函数（拓展考点）。
\end{enumerate}
\noindent\textbf{补充：}建模两类方法：分析法（已知机理，物理定律）、实验辨识法（机理未知，实测输入输出）；同一物理系统可多模型、不同物理系统可共用同一数学模型（例：RLC二阶电路=弹簧-质量-阻尼二阶系统）。

\medskip

\section{控制系统的时域数学模型——微分方程}
\subsection{微分方程的基本形式}

\paragraph{通用非线性微分方程（全系统通用）}
\[
\begin{split}
a_n \frac{\mathrm{d}^n f_n(y(t))}{\mathrm{d}t^n} &+ a_{n-1} \frac{\mathrm{d}^{n-1} f_{n-1}(y(t))}{\mathrm{d}t^{n-1}} + \dots + a_1 \frac{\mathrm{d} f_1(y(t))}{\mathrm{d}t} + a_0 f_0(y(t)) \\
&= b_m \frac{\mathrm{d}^m g_m(u(t))}{\mathrm{d}t^m} + b_{m-1} \frac{\mathrm{d}^{m-1} g_{m-1}(u(t))}{\mathrm{d}t^{m-1}} + \dots + b_1 \frac{\mathrm{d} g_1(u(t))}{\mathrm{d}t} + b_0 g_0(u(t))
\end{split}
\]
其中 $a_i,b_i$ 为系数，$y(t)$输出，$u(t)$输入，$n\ge m$。

\medskip

\paragraph{线性定常(LTI)标准微分方程（核心公式）}
\[
\begin{split}
a_n \frac{\mathrm{d}^n y(t)}{\mathrm{d}t^n} &+ a_{n-1} \frac{\mathrm{d}^{n-1} y(t)}{\mathrm{d}t^{n-1}} + \dots + a_1 \frac{\mathrm{d}y(t)}{\mathrm{d}t} + a_0 y(t) \\
&= b_m \frac{\mathrm{d}^m u(t)}{\mathrm{d}t^m} + b_{m-1} \frac{\mathrm{d}^{m-1} u(t)}{\mathrm{d}t^{m-1}} + \dots + b_1 \frac{\mathrm{d}u(t)}{\mathrm{d}t} + b_0 u(t)
\end{split}
\]

\medskip

\paragraph{系统分类与判定规则}
\begin{itemize}
    \item \textbf{线性系统：}输入输出全为一次项；\textbf{非线性：}含$y^2,y\dot y$等非线性项
    \item \textbf{定常系统：}全部系数$a_i,b_i$为常数；\textbf{时变：}任意系数含$t$
    \item LTI示例：$\dot y+ay=u$；LTV示例：$\dot y+a(t)y=bu$
    \item NLTI示例：$\dot y-y^2=bu$；NLTV示例：$\dot y-y^2=b(t)u$
\end{itemize}

\subsection{建立微分方程（建模三步必考）}
\begin{enumerate}
    \item 列元件物理方程；\item 消去中间变量；\item 输出导数放左侧、输入导数放右侧，降幂整理。
\end{enumerate}

\paragraph{电路建模（KVL/KCL）}
RLC串联：输入$u_i$、输出$u_o=u_C$
\[
LC\frac{\mathrm{d}^2 u_o}{\mathrm{d}t^2}+RC\frac{\mathrm{d}u_o}{\mathrm{d}t}+u_o=u_i(t)
\]

\medskip

\paragraph{机械元件与建模}
\begin{table}[htbp]
\centering
\caption{机械元件基本方程}
\begin{tabular}{ll}
\toprule
元件 & 方程 \\
\midrule
质量 & $F=m\ddot x$ \\
弹簧 & $F=k(x_1-x_2)$ \\
阻尼 & $F=f(\dot x_1-\dot x_2)$ \\
转动惯量 & $T=J\ddot\theta$ \\
\bottomrule
\end{tabular}
\end{table}

弹簧-质量-阻尼（外力$F$，位移$y$）：
\[
m\ddot y+f\dot y+ky=F(t)
\]
\noindent\textbf{机电相似：} RLC与弹簧阻尼质量同型二阶方程。

\subsection{线性系统的性质}
\[
x=ax_1+bx_2 \quad\Longrightarrow\quad y=ay_1+by_2
\]

\subsection{非线性系统的线性化（小信号法，必考）}
平衡点$x_0$泰勒展开，略去高阶小量：
\[
\Delta y=\left.\frac{\mathrm{d}f}{\mathrm{d}x}\right|_{x_0}\cdot\Delta x,\quad K=\left.\frac{\mathrm{d}f}{\mathrm{d}x}\right|_{x_0}
\]
\textbf{不可线性化：}继电器、死区、间隙等本质非线性。

\bigskip

\section{控制系统的复域数学模型——传递函数}
\subsection{拉普拉斯变换定义}
\[
F(s)=\mathcal{L}[f(t)]=\int_0^\infty f(t)e^{-st}\mathrm{d}t
\]

\paragraph{常用变换对（必背）}
\begin{table}[htbp]
\centering
\caption{常见拉普拉斯变换对}
\begin{tabular}{ccc}
\toprule
$f(t)$ & $F(s)$ & 条件 \\
\midrule
$\delta(t)$ & $1$ & \\
$1(t)$ & $\dfrac1s$ & \\
$t$ & $\dfrac1{s^2}$ & \\
$\dfrac{t^{n-1}}{(n-1)!}$ & $\dfrac1{s^n}$ & $n=1,2,\dots$ \\
$e^{-at}$ & $\dfrac1{s+a}$ & \\
$\sin\omega t$ & $\dfrac{\omega}{s^2+\omega^2}$ & \\
$\cos\omega t$ & $\dfrac{s}{s^2+\omega^2}$ & \\
$\dfrac{t^{n-1}}{(n-1)!}e^{-at}$ & $\dfrac1{(s+a)^n}$ & $n=1,2,\dots$ \\
\bottomrule
\end{tabular}
\end{table}

\paragraph{拉普拉斯变换性质}
\begin{enumerate}
\item \textbf{线性性质：}$\mathcal{L}[a f_1(t)+b f_2(t)]=aF_1(s)+bF_2(s)$

\item \textbf{微分性质：}
一阶：$\displaystyle \mathcal{L}\left[\frac{\mathrm{d}f(t)}{\mathrm{d}t}\right]=sF(s)-f(0)$
$n$阶：$\displaystyle \mathcal{L}\left[\frac{\mathrm{d}^n f(t)}{\mathrm{d}t^n}\right]=s^nF(s)-s^{n-1}f(0)-\dots-f^{(n-1)}(0)$
零初值：$\mathcal{L}[f^{(n)}(t)]=s^nF(s)$

\item \textbf{积分性质：}
单次积分：$\displaystyle \mathcal{L}\left[\int f(t)\mathrm{d}t\right]=\frac1sF(s)+\frac1s f(0)$
$n$重积分零初值：$\displaystyle \mathcal{L}\left[\underbrace{\int\dots\int}_{n} f(t)\mathrm{d}t^n\right]=\frac1{s^n}F(s)$

\item \textbf{初值定理：}$\displaystyle \lim_{t\to0^+}f(t)=\lim_{s\to\infty}sF(s)$

\item \textbf{终值定理：}$\displaystyle \lim_{t\to\infty}f(t)=\lim_{s\to0}sF(s)$（$sF(s)$极点全在左半平面）

\item \textbf{位移定理：}
时域延迟：$\mathcal{L}\big[f(t-\tau_0)\big]=e^{-\tau_0 s}F(s)$
复频位移：$\mathcal{L}\big[e^{-at}f(t)\big]=F(s+a)$
\end{enumerate}

其他常用性质：
\[
\begin{aligned}
&\mathcal{L}\big[f(t-a)\cdot 1(t-a)\big]=e^{-as}F(s)\quad(a>0),\\
&\mathcal{L}\big[t f(t)\big]=-\frac{\mathrm{d}F(s)}{\mathrm{d}s},\\
&\mathcal{L}\big[\tfrac1t f(t)\big]=\int_s^\infty F(s)\mathrm{d}s,\\
&\mathcal{L}\big[f(\tfrac ta)\big]=aF(as)\quad(a>0).
\end{aligned}
\]

\medskip

\paragraph{留数定理（拉普拉斯逆变换）}
若$F(s)$为有理分式，极点$s=p_i\;(i=1,2,\dots,k)$，则
\[
f(t)=\mathcal{L}^{-1}[F(s)]=\frac1{2\pi j}\int_{\sigma-j\infty}^{\sigma+j\infty}F(s)e^{st}\mathrm{d}s=\sum_{i=1}^k \mathrm{Res}\big[F(s)e^{st},\,s=p_i\big],\quad t>0
\]
\begin{itemize}
    \item \textbf{一阶极点 $s=p$：}$\displaystyle \mathrm{Res}\big[F(s)e^{st},p\big]=\lim_{s\to p}(s-p)F(s)e^{st}$
    \item \textbf{$r$阶重极点 $s=p$：}$\displaystyle \mathrm{Res}=\frac1{(r-1)!}\lim_{s\to p}\frac{\mathrm{d}^{r-1}}{\mathrm{d}s^{r-1}}\big[(s-p)^r F(s)e^{st}\big]$
\end{itemize}

\subsection{传递函数的概念与说明}
\paragraph{定义（零初始条件下）}
系统输出量的拉氏变换与输入量的拉氏变换之比：
\[
G(s) = \frac{Y(s)}{R(s)} = \frac{b_m s^m + b_{m-1} s^{m-1} + \dots + b_1 s + b_0}{a_n s^n + a_{n-1} s^{n-1} + \dots + a_1 s + a_0}
\]

\paragraph{传递函数的八点详细说明}
\begin{enumerate}
    \item 它是复域数学模型，和系统微分方程一一对应；
    \item 只取决于系统或元件的结构和参数，与系统输入信号无关；
    \item 形式与输入、输出信号的选取位置密切相关；
    \item 分子分母是$s$的降幂有理多项式，系数均为实数，零点、极点只可为实数或共轭复数；
    \item 物理可实现系统分母阶次大于或等于分子阶次（$n\ge m$）；
    \item 无法直接反映系统及元件的物理结构；
    \item 令分母多项式等于零，可得到系统的特征方程；
    \item 传递函数和系统单位脉冲响应互为拉普拉斯变换对。
\end{enumerate}

\paragraph{传递函数的局限性}
\begin{enumerate}
    \item 原则上不反映非零初始条件时系统响应的全部信息；
    \item 适合用于描述单输入单输出（SISO）系统；
    \item 只能用于表示线性定常系统。
\end{enumerate}

\medskip

\paragraph{零初始条件说明（考试常考点）}
\begin{itemize}
    \item 输入量在 $t\ge0$ 作用于系统，$t=0^-$ 时输入及其各阶导数均为零；
    \item 输入加入前系统处于稳定工作状态，$t=0^-$ 时输出及其各阶导数也为零。
\end{itemize}

\medskip

\paragraph{由微分方程推导传递函数}
对线性定常系统微分方程：
\[
a_n y^{(n)}(t) + a_{n-1} y^{(n-1)}(t) + \dots + a_1 \dot{y}(t) + a_0 y(t) = b_m r^{(m)}(t) + \dots + b_1 \dot{r}(t) + b_0 r(t)
\]
零初始条件下，取拉普拉斯变换：
\[
(a_n s^n + a_{n-1} s^{n-1} + \dots + a_1 s + a_0)Y(s) = (b_m s^m + \dots + b_1 s + b_0)R(s)
\]
整理即得传递函数表达式。

\medskip

\paragraph{传递函数的三种标准形式（必考）}
\subparagraph{1. 零、极点形式（根轨迹分析常用）}
将传递函数因式分解，写成零极点乘积形式：
\[
G(s) = \frac{Y(s)}{R(s)} = \frac{b_m s^m + b_{m-1} s^{m-1} + \dots + b_1 s + b_0}{a_n s^n + a_{n-1} s^{n-1} + \dots + a_1 s + a_0}
\]
\[
G(s) = \frac{b_0}{a_0} \cdot \frac{(s-z_1)(s-z_2)\cdots(s-z_m)}{(s-p_1)(s-p_2)\cdots(s-p_n)} = K^* \cdot \frac{\prod_{i=1}^m (s-z_i)}{\prod_{j=1}^n (s-p_j)}
\]
其中：
\begin{itemize}
    \item 零点：$z_i\;(i=1,2,\dots,m)$，使分子为零的复频率；
    \item 极点：$p_j\;(j=1,2,\dots,n)$，使分母为零的复频率，也是微分方程的特征根；
    \item 根轨迹增益（传递系数）：$K^* = \dfrac{b_0}{a_0}$。
\end{itemize}

\subparagraph{2. 时间常数形式（频率法分析常用）}
将传递函数整理为以“时间常数”为系数的标准形式，方便绘制频率特性：
\[
G(s) = \frac{b_m}{a_n} \cdot \frac{(\tau_1 s+1)(\tau_2^2 s^2 + 2\zeta\tau_2 s + 1)\cdots(\tau_i s+1)}{(T_1 s+1)(T_2^2 s^2 + 2\zeta T_2 s + 1)\cdots(T_j s+1)}
\]
其中：
\begin{itemize}
    \item 一次因子 $(\tau_i s+1)$ 对应实数零极点；
    \item 二次因子 $(T_j^2 s^2 + 2\zeta T_j s + 1)$ 对应共轭复数零极点；
    \item 时间常数：$\tau_i$（分子）和 $T_j$（分母）；
    \item 传递系数（静态增益）：$K = \dfrac{b_m}{a_n}$。
\end{itemize}

\subparagraph{3. 两种形式的转换关系}
从零极点形式到时间常数形式的转换，需要将零极点因子改写为“一次项+常数”的形式：
\[
G(s) = K^* \cdot \frac{\prod_{i=1}^m (s-z_i)}{\prod_{j=1}^n (s-p_j)} = K^* \cdot \frac{\left( \prod_{i=1}^m \frac{1}{-z_i} \right) \prod_{i=1}^m \left( \frac{1}{-z_i} s + 1 \right)}{\left( \prod_{j=1}^n \frac{1}{-p_j} \right) \prod_{j=1}^n \left( \frac{1}{-p_j} s + 1 \right)}
\]
两种增益的关系为：
\[
K = K^* \cdot \frac{\prod_{i=1}^m (-z_i)}{\prod_{j=1}^n (-p_j)}
\]

\subsection{典型环节的传递函数（8类全细化，必考）}
控制系统复杂传递函数都可拆解为以下基础环节的组合：
\begin{enumerate}
    \item \textbf{放大环节（比例环节）}
    时域关系：$y(t)=Kr(t)$，传递函数：$\displaystyle G(s)=\frac{Y(s)}{R(s)}=K$，几乎所有控制系统都包含该环节。

    \item \textbf{惯性环节}
    微分方程：$\displaystyle T\frac{\mathrm{d}y(t)}{\mathrm{d}t}+y(t)=r(t)$，传递函数：$\displaystyle G(s)=\frac{Y(s)}{R(s)}=\frac1{Ts+1}$，$T$为时间常数；$T=0$时退化为放大环节。

    \item \textbf{积分环节}
    微分方程：$\displaystyle y(t)=\int r(t)\mathrm{d}t$，传递函数：$\displaystyle G(s)=\frac{Y(s)}{R(s)}=\frac1s$，输出为输入对时间的积分。

    \item \textbf{震荡环节}（$0\le\zeta<1$）
    微分方程：$\displaystyle T^2\frac{\mathrm{d}^2y(t)}{\mathrm{d}t^2}+2\zeta T\frac{\mathrm{d}y(t)}{\mathrm{d}t}+y(t)=r(t)$，
    传递函数：$\displaystyle G(s)=\frac{Y(s)}{R(s)}=\frac{\omega_n^2}{s^2+2\zeta\omega_n s+\omega_n^2}$。

    \item \textbf{纯微分环节}
    微分方程：$\displaystyle y(t)=\frac{\mathrm{d}r(t)}{\mathrm{d}t}$，传递函数：$\displaystyle G(s)=\frac{Y(s)}{R(s)}=s$，输出为输入对时间的微分。

    \item \textbf{一阶微分环节}
    微分方程：$\displaystyle y(t)=\tau\frac{\mathrm{d}r(t)}{\mathrm{d}t}+r(t)$，传递函数：$\displaystyle G(s)=\frac{Y(s)}{R(s)}=\tau s+1$。

    \item \textbf{二阶微分环节}
    微分方程：$\displaystyle y(t)=\tau^2\frac{\mathrm{d}^2r(t)}{\mathrm{d}t^2}+2\zeta\tau\frac{\mathrm{d}r(t)}{\mathrm{d}t}+r(t)$，
    传递函数：$\displaystyle G(s)=\frac{Y(s)}{R(s)}=\tau^2 s^2+2\zeta\tau s+1$。

    \item \textbf{延迟环节}
    微分方程：$y(t)=r(t-\tau)$，传递函数：$\displaystyle G(s)=\frac{Y(s)}{R(s)}=e^{-\tau s}$，$\tau$为延迟时长。
\end{enumerate}

\bigskip

\section{控制系统的方框图和信号流图}
\subsection{2.3.1 方框图的概念和绘制}
方框图是用图形直观描述系统各元件信号传递关系的数学模型，能清晰展示信号流向、环节连接方式与系统结构，方便求解复杂系统传递函数。

对于结构复杂的多回路系统，单纯利用方框图化简方法求取系统整体传递函数，计算过程会变得十分繁琐困难。

\subsection{2.3.2 方框图化简}
\begin{enumerate}
    \item \textbf{串联环节等效}：总传递函数 $G(s) = G_1(s)G_2(s)$；
    \item \textbf{并联环节等效}：总传递函数 $G(s) = G_1(s) + G_2(s)$；
    \item \textbf{反馈闭环等效}：负反馈 $\displaystyle \Phi(s) = \frac{G(s)}{1 + G(s)H(s)}$，正反馈分母改为 $1-G(s)H(s)$。
\end{enumerate}

\subsection{2.3.3 信号流图的概念和绘制}
信号流图是表示系统中各变量间相互关系以及信号传递的另一种图解方法。
\begin{itemize}
    \item 节点：表示系统的变量；
    \item 支路：相当于乘法器，信号流经支路时，被乘以支路增益而变换为另一信号；
    \item 信号在支路上只能沿箭头单向传递；
    \item 系统节点变量的设置不唯一，因此同一个系统的信号流图不是唯一的。
\end{itemize}

\paragraph{节点运算规则}
节点同时具备两种作用：
\begin{enumerate}
    \item 对所有来自流入支路的信号作加法运算；
    \item 将流入信号之和，传输给所有流出支路。
\end{enumerate}
\[
\begin{cases}
w = au+bv \\
x = cw = c(au+bv) \\
y = dw = d(au+bv)
\end{cases}
\]

\paragraph{信号流图核心术语}
\begin{itemize}
    \item \textbf{源节点（输入节点）}：只有信号输出，无信号输入；
    \item \textbf{阱节点（输出节点）}：只有信号输入，无信号输出；
    \item \textbf{混合节点}：既有输入支路，又有输出支路；
    \item \textbf{前向通路}：信号从输入节点至输出节点传递，每个节点只通过1次的通路；
    \item \textbf{单独回路}：起点和终点为同一节点，信号遍历节点不重复的闭合通路；
    \item \textbf{回路增益}：回路中所有支路增益的乘积；
    \item \textbf{不接触回路}：多个回路之间没有任何公共节点，反之称为接触回路。
\end{itemize}

\paragraph{方框图 $\leftrightarrow$ 信号流图对应转换规则}
方框图单元 $\rightarrow$ 信号流图单元：
输入相加点 $\rightarrow$ 混合节点；引出点 $\rightarrow$ 混合节点；方框环节 $\rightarrow$ 增益支路；信号连线 $\rightarrow$ 有向支路。

\subsection{2.3.4 梅森增益公式}
复杂系统无需化简框图，可直接通过梅森公式求解系统总传递函数，公式为：
\[
P = \frac{1}{\Delta} \sum_{k=1}^n P_k \Delta_k
\]

\paragraph{参数定义}
\begin{itemize}
    \item $P$：系统总传输增益（闭环总传递函数）；
    \item $P_k$：第 $k$ 条前向通路的通路总增益；
    \item $\Delta$：信号流图特征式
    \[
    \Delta = 1 - \sum L_a + \sum L_b L_c - \sum L_d L_e L_f + \dots
    \]
    \begin{itemize}
        \item $\sum L_a$：所有单独回路增益之和；
        \item $\sum L_b L_c$：每两个互不接触回路增益乘积之和；
        \item $\sum L_d L_e L_f$：每三个互不接触回路增益乘积之和；
    \end{itemize}
    \item $\Delta_k$：第 $k$ 条前向通路的余子式
    \begin{itemize}
        \item 若通路与所有回路都接触：$\boldsymbol{\Delta_k=1}$；
        \item 若通路与所有回路都不接触：$\boldsymbol{\Delta_k=\Delta}$；
    \end{itemize}
\end{itemize}

\bigskip

\section{第2章 小结}
\subsection{经典控制理论的数学模型体系}
经典控制理论中常用的控制系统数学模型包含四类：
\begin{itemize}
    \item 微分方程（时域最原始模型）
    \item 传递函数（复域核心模型）
    \item 方框图（图形化模块模型）
    \item 信号流图（简化图形化模型，适配梅森公式）
\end{itemize}

\subsection{系统分析与综合的闭环逻辑}
由数学模型得到系统模型后，可开展双向工作：
\begin{enumerate}
    \item \textbf{系统分析（正向推导）}：由系统模型计算控制系统的各项性能指标，判断其是否满足设计要求，分为三类方法：
    \begin{itemize}
        \item 时域法
        \item 复域法
        \item 频域法
    \end{itemize}
    \item \textbf{系统综合（反向设计）}：根据性能指标要求，寻求改进系统性能、使其满足设计规范的校正方法，形成闭环优化流程。
\end{enumerate}

% ===== 恢复全局默认间距（避免影响后续章节）=====
\setlength{\parskip}{0pt plus 1pt}
\setlist{leftmargin=2em, itemsep=0.2em, parsep=0pt}
\renewcommand{\arraystretch}{1}
\setlength{\tabcolsep}{6pt}
% ============================================